\section{Система обучения}

Фреймворк обучения рассматривается практически такой же, как был в работе \cite{planes}.
Мы рассматриваем фреймворк с алгоритмом обучения Deep-Q Network и $\epsilon$-жадной политикой.

\subsection{DQN-алгоритм}

Алгоритм Deep-Q Network - это стандартный безмодельный алгоритм обучения агента в виде 
нейронной сети на принятия решения в среде.

Каждое взаимодействие между агентом или 
несколькими агентами с окружением хранится в виде элемента Марковской цепи решений (MDP), состоящего из:

\begin{itemize}
\item Некоторого состояния окружения \textbf{$s_0$}
\item Набора действий каждого агента при этом состоянии (одно на каждого агента)
\item Награда за совершение этих действий
\item Нового состояния-результата \textbf{$s_1$}
\end{itemize}

Элементы (или переходы Марковской цепи) хранятся в \textbf{воспроизводящей памяти} $RM$
ограниченного объёма $N$ - она применяется в обучении нейронной сети агента.

В алгоритме используется обучение двух нейронных сетей:
\begin{itemize}
\item Тренировочная сеть $Q_{policy}$ - обучается явным образом на каждом шаге эпизода
\item Целевая сеть $Q_{target}$ - обучается неявно путём линейного смещения параметров в 
    сторону параметров 
	тренировочной сети с небольшим коэффициентом, тем самым ограничивая её от переобучения:
	\begin{equation}
	\label{eq:tar_pol}
	Q'_{target} = (1-\tau)Q_{target} + \tau Q_{policy}
	\end{equation}
\end{itemize}

Именно целевая нейронная сеть является результирующей политикой агента после обучения.
Тренировочная сеть используется для основной коммуникации с окружением и генерации MDP-элементов,
заносимых в воспроизводящую память для обучения.

Для выбора действия, нейронная сеть агента предсказывает награду, которую он получит за совершение $j$-го действия.
В итоге окружением совершается действие, имеющее максимальную предсказанную награду.

\begin{align*}
	\textbf{$\hat r_t$} &= Q(\textbf{$s_t$}) \\
	r_t(a^j) &= \hat r^j
	a^*_t &= argmax_a (r_t(a))
\end{align*}
 
Обучение тренировочной сети проходит после каждого ввода нового MDP-элемента со следующим алгоритмом:
\begin{enumerate}
\item Берётся сет из $N_B$ случайно взятых переходов из воспроизводящей памяти
\item Расчитывается целевая награда сета - сумма награды после изначального состояния и награды 
	за действия выбранные целевой политикой на новом состоянии $s_{t+1}$:
	\begin{equation}
	\label{eq:r_tar}
	r_{target} = r_t + \gamma \tau Q_{target}(s_{t+1})
	\end{equation}
	, где $\gamma$ - это коэффициент дисконтирования награды, снижающий значимость 
	награды, полученную на более поздних ходах 
\item Предсказанная награда сета - это награда за действия выбранные тренировочной сетью 
		на изначальном состоянии:
	\begin{equation}
	\label{eq:r_pol}
	r_{policy} = Q_{policy}(s_{t+1})
	\end{equation}	
\item 
	Обучение тренировочной сети происходит путём оптимизации функции потерь Хубера между целевой
	\eqref{eq:r_tar} и предсказанной \eqref{eq:r_pol} наградами.
	\begin{equation}
	Q'_{policy} = Q_{policy} + \beta \nabla_Q L_{\delta}(r_{target}, r_{policy})
	\end{equation}
	Тем самым, тренировочная модель "мотивирована" искать шаги, позволяющие за один ход получить 
	награду, ранее достигаемую за два хода целевого агента, и следовательно - достигать цели 
	наиболее быстрыми методами.
\item Наконец, производится обновление параметров целовой политики по уравнению \eqref{eq:tar_pol}.
\end{enumerate}

\subsection{Уровень исследования}

Уровень исследования - это вероятность, с которой на каждом шаге генрации MDP-перехода 
обучающийся агент выберет случайное действие, вместо результата собственной обработки состояния 
окружения. 
Она является функцией от номера эпизода и рассчитывается как:

\[ \epsilon(n) = 0.05 + 0.9\ max(1 - \frac{n+n_{r}}{N_{total}}, 0),\ n \leq N_{total}\]
где $N_{total}$ — общее количество обучающих эпизодов, а $n_{r}$ — количество последних эпизодов,
в которых скорость исследования минимальна.

В случае нашей основной задачи окружения, в виду существования обширной запрещенной
зоны, случайный выбор действия часто приводит к большому числу провальных попыток во время обучения.

Однако данная функция всё же сохраняет особую важность в глубоком обучении с подкреплением, 
так как придаёт памяти разнообразие состояний окружений и действий агента в обучающем "датасете".

В работе $n_{r}$ фиксируется как в процессе обучения, так и в процессе послешокового 
переобучения и составляет 500 эпизодов.